

The report-resolved continuation theorem is exact for linear marked task equalities of a supplied time-homogeneous finite-return routine. It does not identify the reachable positive cone with every positive matrix on its support, nor replace a one-sided reachable cost bound by an unsupported operator inequality. The reset-instrument classification applies to failure branches which physically prepare known states and to first-success commitment. It includes arbitrary finite rank-one families, but not arbitrary quantum instruments or an arbitrary-rank projection without input-erasing failure. The zero-risk spectrum is exact for the acquired profile; its physical threshold has the same converse and implementation hypotheses as the exact transfer theorem. These distinctions separate the new structural results from the independent wider research goals.


The primitive finite-model route in Section~\ref{sec:primitive-risk} removes the reset and reference assumptions on its own declared interface. It supplies a common fallback, monotone two-sided approximations to the least feasible cost, and bounded-risk approximations at every strictly interior feasible budget with counted finite implementations. It does not replace the preparation-free all-initial-state quantifier by an arbitrary restricted preparation set, and it assumes finite models, finite globally legal commands/reports, time homogeneity and bounded positive call costs. It is not unknown-kernel statistical learning, a universal closed formula for the comparison profile, or optimal computational complexity. The six-regime diagram solves one explicit mixed-strata comparator with its original asymptotic qualification intact.
